\begin{figure}[tb]
\centering
\begin{tikzpicture}[x=1.55cm,y=1.55cm,>=Latex,font=\small]
 \begin{scope}
  \clip (-.25,-1.1) rectangle (4.4,1.65);
  \fill[blue!5] (0,0) circle (3);
 \end{scope}
 \fill (0,0) circle (1.4pt) node[left] {$0$};
 \draw[->,gray] (0,0)--(3,0) node[midway,below] {$r$};
 \node[blue!45!black] at (1.15,1.34) {base ball $B_r$};
 \draw[red!65!black,dashed,thick] (3.08,.65) circle (.85);
 \fill[red!65!black] (3.08,.65) circle (1.4pt);
 \draw[red!65!black] (3.08,.65)--(3.89,.91)
      node[midway,above] {$\ell$};
 \node[red!65!black,anchor=west] at (3.57,1.51) {failure ball};
 \draw[red!65!black,->] (3.6,1.39)--(3.57,1.25);
 \draw[gray,fill=gray!12] (2.63,.3) rectangle (3.18,.9);
 \begin{scope}
  \clip (-.25,-1.1) rectangle (4.4,1.65);
  \draw[blue!65!black,thick] (0,0) circle (3);
 \end{scope}
 \node[blue!65!black,anchor=west] at (3.02,-.33) {$\rho=r$};
 \fill (2.78,.58) circle (1.4pt) node[above left] {$y$};
 \fill[green!14] (1.73,.58) circle (.32);
 \draw[green!40!black,thick] (1.73,.58) circle (.32);
 \draw[->,thick,orange!75!black] (2.58,.58)--(2.12,.58);
 \node[green!35!black,align=center,anchor=north] at (1.45,-.32)
      {inward subball\\with positive height mass};
 \draw[green!35!black,->] (1.64,-.29)--(1.73,.2);
 \node[anchor=north,align=center] at (3.81,-.71)
      {cell crossing\\the kernel boundary};
 \draw[gray,->] (3.66,-.65)--(3.14,.3);
\end{tikzpicture}
\caption{An inward subball controls the boundary weights on the reference center. A failure at
scale $\ell$ supplies separated projection entries on every sufficiently small comparable
subball in its comparison region. Moving such a subball inward by $O(\ell)$
gives positive height mass where the polynomial kernel has a controlled lower
bound. This pays for a local cell even when that cell crosses $\rho=r$. The
picture is schematic; no circular shape of a failure region is assumed.}
\label{fig:collar}
\end{figure}
